\section{Cost management}
\label{sec:cost}

The cost baseline is the WBS in Figure~\ref{fig:wbs} priced at Level~3. Seventeen control accounts sum to a work estimate of \$1{,}500{,}000. Contingency of \$78{,}890 is sized from residual expected monetary value plus four named residual items, not from a round-number markup, so the authorised budget at completion is \$1{,}578{,}890.

\subsection{Estimating basis and assumptions}
\label{sec:cost-basis}

The estimate is bottom-up at Level~3, with analogous plant and commodity rates beside the hours. Loaded 2026 Australian-dollar rates are used throughout: systems-engineering and project-management time at \$150/h, pad technicians at \$95/h, cryogenic and range-safety specialists at \$165/h, and surge technicians at \$100/h. The Manufacturing Award C10 ordinary hourly rate is \$29.45 from the first full pay period on or after 1~July~2026 \citep{fwc2026}. Charge-out is C10 multiplied by a named employment on-cost, then utilisation and overhead---the product is in the notes to Table~\ref{tab:wbs-cost}. Those factors do not change the \$1{,}500{,}000 lines. The 100~t crawler for 23\,m stacking is priced at \$4{,}200 per wet day, including operator, from published Queensland plant-hire schedules cross-checked against Rawlinsons \citep{rawlinsons2025}.

Inclusions are the seventeen packages in Figure~\ref{fig:wbs}: mobilisation through demobilisation at Bowen. Factory manufacture, engine research and development, and freight to the receiving gate remain out of scope, so A-100 is a zero-dollar external driver rather than a cost account. The four-person surge crew (\$22{,}400) is inside WBS~1.3.3 as planned capacity. The \$21{,}000 hybrid compression cost in Section~\ref{sec:schedule-compression} is a contingent draw against the \$78{,}890 reserve if R-14 is realised; it is not in the \$1{,}500{,}000 performance measurement baseline.

The 7~August charter does not authorise field spend. Permitting and Juru work start on 10~August because those durations must finish before hot fire. Other August--September pad cost is a PEP choice: planned value already on the clock by Gate~1 (A-111, 2~October) is \$602{,}301 (40.2\% of the work PMB), including stacking in progress. Gate~1 locks remaining spend; it does not pretend the campaign is only a mobilisation paper.

\subsection{WBS cost baseline}
\label{sec:cost-wbs}

Table~\ref{tab:wbs-cost} is the rollup. Pad logistics (1.2) is \$518{,}000 (34.5\% of the work estimate), pad testing (1.3) \$414{,}000 (27.6\%), governance (1.1) \$300{,}000 (20.0\%), and flight/closeout (1.4) \$268{,}000 (17.9\%).

\begin{table}[!htbp]
\centering
\footnotesize
\setlength{\tabcolsep}{4.2pt}
\caption[Bottom-up cost baseline by WBS]{Bottom-up cost baseline by Level-3 WBS package. Totals equal the work estimate of \$1{,}500{,}000. A-100 is omitted because it is not a cost account.}
\label{tab:wbs-cost}
\begin{tabularx}{\textwidth}{@{}l p{4.15cm} r X@{}}
\toprule
WBS & Package & Amount (AUD) & Basis \\
\midrule
1.1.1 & PEP and integration & 72{,}000 & 480~h PM @ \$150 \\
1.1.2 & Governance and EVM & 48{,}000 & 320~h planner/cost @ \$150 \\
1.1.3 & ASA / CASA / GBRMPA & 88{,}000 & ASA \$42k $+$ CASA \$28k $+$ GBRMPA \$18k \\
1.1.4 & Juru monitors and heritage & 92{,}000 & Two monitors + survey \\
1.2.1 & Receiving and cleanroom & 68{,}000 & Technician hours + consumables \\
1.2.2 & Crane and 23\,m stacking & 218{,}000 & 8 lift + 16 weld/NDT/bolt-up d $\times$ \$4{,}200 $=$ \$100{,}800 + riggers \\
1.2.3 & Cryo GSE & 180{,}000 & Valves/piping + specialist crew \\
1.2.4 & RF array calibration & 52{,}000 & Crew + analyser hire \\
1.3.1 & Avionics power-on & 58{,}000 & Engineers @ \$150/h + EGSE \\
1.3.2 & WDR (cold) & 160{,}000 & Includes bulk LOX 25~t and LN2 \$68{,}500 \\
1.3.3 & Static fire & 148{,}000 & Includes surge hire \$22{,}400 \\
1.3.4 & Post-fire / pre-kitted spare & 48{,}000 & Engineers + spare handling \\
1.4.1 & LRR and RSO & 38{,}000 & Review board; payload witnesses \\
1.4.2 & Range exclusion / marine & 58{,}000 & Patrol boats + NOTAM \\
1.4.3 & Launch and telemetry capture & 82{,}000 & Launch-day operations \\
1.4.4 & Decrypt and handover & 24{,}000 & Five-day engineering handover \\
1.4.5 & Demobilisation and audit & 66{,}000 & Site restore + audit \\
\midrule
\textbf{Base} & & \textbf{1{,}500{,}000} & \mbox{} \\
\bottomrule
\end{tabularx}

\vspace{0.35em}
{\footnotesize\textit{Notes.} Labour: MA000010 C10 ordinary \$29.45/h (ppc 1~July~2026) from Annual Wage Review 2026 \citep{fwc2026}. Employment on-cost $=$ super 12\% $+$ annual leave 7.69\% $+$ personal leave 3.85\% $+$ public holidays 4.23\% $+$ LSL 2.17\% $+$ 17.5\% leave loading on annual leave $+$ workers compensation (1.8\% professional / 4.5\% pad) $+$ QLD payroll tax 4.75\% ($1.378\times$ professional, $1.405\times$ pad). Charge-out then divides by utilisation and applies overhead, with no consultancy fee: engineer/PM $\$29.45 \times 1.378 / 0.72 \times 2.66 = \$150$; technician $\$29.45 \times 1.405 / 0.85 \times 1.95 = \$95$; specialist $\$29.45 \times 1.378 / 0.60 \times 2.44 = \$165$. Surge is C10 $\times$ 1.5 overtime $\times$ pad on-cost $/ 0.85 \times 1.37 = \$100$. Crane: 8 lift days and 16 weld/NDT/bolt-up days at the same wet \$4{,}200 (QLD 100~t crawler band about \$3{,}800--\$4{,}500/day), cross-checked against Rawlinsons \citep{rawlinsons2025}. Weather dollars sit only in RES (4 extra wet days). 1.1.3: ASA 280~h @ \$150 $=$ \$42{,}000; CASA CRIS airspace instrument \$28{,}000; GBRMPA permission \$18{,}000 \citep{asa2025,casa2024,gbrmpa2019}. LOX/LN2 \$68{,}500 is a bulk industrial-gas analogue inside 1.3.2. Surge \$22{,}400 $=$ $4 \times 7$~d $\times$ 8~h $\times$ \$100 and sits in 1.3.3, not in contingency.}
\end{table}

The largest package is crane mobilisation and stacking (1.2.2, \$218{,}000) because A-123 occupies 24 calendar days on the launch-critical path. The wet crawler is on hire for that whole window: 8 lift days and 16 weld, NDT and bolt-up days at \$4{,}200, which is \$100{,}800, with the balance in rigger labour. Those 16 days are plant duration, not weather. Weather dollars sit in one place only: RES four extra wet days at \$4{,}200 $=$ \$16{,}800. R-06 is a schedule control (dawn-lift window) with residual EMV \$0, so it is not a second crane line. Wet Dress Rehearsal (1.3.2) carries \$68{,}500 of bulk liquid oxygen and nitrogen inside \$160{,}000. Static fire (1.3.3) carries the planned surge inside \$148{,}000. The seventeen lines sum to \$1{,}500{,}000 with no miscellaneous remainder. Package 1.2.1 remains one control account for A-121 and A-122.

\subsection{Contingency}
\label{sec:cost-contingency}

Contingency is not a 15\% story. Table~\ref{tab:contingency} adds a residual expected-monetary-value pool to four named residual items. Treatments already sit in the \$1{,}500{,}000 base (surge hire, dawn-lift window, dual-path recorders, pre-kitted spare), so the pool uses residual $P'_{\%}\times I_{\$}$, not inherent $P_{\%}$. Summing the ten cost-bearing risks R-01, R-02, R-03, R-04, R-05, R-07, R-09, R-10, R-13 and R-14 gives \$35{,}760. The residual-uncertainty (RES) allowance is itemised in the same table and sums to \$43{,}130. Vehicle freight is out of scope and is not in RES. The two components sum to \$78{,}890, which is 5.26\% of \$1{,}500{,}000, and the authorised $BAC$ is therefore \$1{,}578{,}890.

\begin{table}[H]
\centering
\small
\caption[Contingency build-up]{Contingency build-up. The 5.26\% ratio is residual EMV plus four named residual items, not an input markup.}
\label{tab:contingency}
\begin{tabular}{@{}p{8.4cm} r p{4.1cm}@{}}
\toprule
Component & Amount (AUD) & Rule \\
\midrule
$\Sigma$ residual EMV of ten cost-bearing risks & 35{,}760 & Residual $P'_{\%}\times I_{\$}$; treatments already in the base \\
LOX/LN2 price variance (18\% of \$68{,}500 gases) & 12{,}330 & In-scope commodity; not freight \\
Crane weather (only weather \$): 4~d $\times$ \$4{,}200 & 16{,}800 & Extra wet days; R-06 is schedule-only \\
GN2/N2 purge: 40 G-size cylinders $\times$ \$200 & 8{,}000 & Named multiply-out \\
Cryogenic kit: 3 $\times$ 3~m LOX hoses $\times$ \$1{,}400 $+$ seal kit \$1{,}800 & 6{,}000 & Named multiply-out \\
\midrule
\textbf{RES subtotal} & \textbf{43{,}130} & Itemised; not a plug to 15\% \\
\textbf{Contingency reserve} & \textbf{78{,}890} & 5.26\% of base \\
\textbf{$BAC$} & \textbf{1{,}578{,}890} & Base + contingency \\
\bottomrule
\end{tabular}
\end{table}

Several register rows are kept out of the \$35{,}760 so that money is not counted twice. R-08 (pad-civil defect) is contract-transferred under the design-and-construct package and carries no Gilmour EMV. R-06, R-11, R-15, R-16, R-17 and R-18 are schedule or HSE exposures without a residual dollar draw. R-12 is the TestFlight~1 in-flight failure mode; the vehicle is factory product, so site-ops EMV is \$0. If R-14 is realised and Option~C is invoked, the \$21{,}000 is drawn from the \$78{,}890 reserve. It is not added to the \$1{,}500{,}000 work baseline. There is no management reserve above $BAC$: unknown unknowns would require a board variation this plan does not authorise.

\subsection{Time-phased budget and cash flow}
\label{sec:cost-cashflow}

The performance measurement baseline is the work estimate spread across the same dates as Table~\ref{tab:precedence}. Table~\ref{tab:cashflow} and Figure~\ref{fig:scurve} use a uniform daily burn on each activity. October is the peak month because stacking, Wet Dress Rehearsal and the five-second static fire all sit there. December is a residual on 1~December, not a parking place for unused contingency.

\begin{table}[H]
\centering
\small
\caption[Time-phased planned value]{Time-phased planned value of the \$1{,}500{,}000 work baseline. Contingency is not monthly spend.}
\label{tab:cashflow}
\begin{tabular}{@{}lrr@{}}
\toprule
Month & Period PV (AUD) & Cumulative PV (AUD) \\
\midrule
Aug 2026 & 155{,}951 & 155{,}951 \\
Sep 2026 & 407{,}188 & 563{,}139 \\
Oct 2026 (peak: stack, WDR, static fire) & 628{,}861 & 1{,}192{,}000 \\
Nov 2026 & 302{,}000 & 1{,}494{,}000 \\
Dec 2026 & 6{,}000 & \textbf{1{,}500{,}000} \\
\bottomrule
\end{tabular}
\end{table}

The S-curve therefore terminates at \$1{,}500{,}000. Figure~\ref{fig:scurve} draws $BAC = \$1{,}578{,}890$ as a horizontal authorisation line. Planned value at M-4 (15~October) is \$865{,}578---the same total used in Section~\ref{sec:delivery-evm}. Planned value at Gate~1 (2~October) is \$602{,}301. Option~C's \$21{,}000 is excluded from every period in Table~\ref{tab:cashflow}. If it is spent, it appears as actual cost charged to contingency, not as a rewrite of the S-curve.

\begin{figure}[H]
\centering
\includegraphics[width=0.98\linewidth]{cash_flow_scurve.pdf}
\caption[Planned-value S-curve]{Cumulative planned-value S-curve for the Bowen campaign. Period bars are monthly PV; the teal line is cumulative PV and ends at the \$1{,}500{,}000 work baseline. The crimson dotted line is $BAC$ (\$1{,}578{,}890) as an authorisation ceiling, not December cash. Contingency remains undistributed. Option~C is not in the curve.}
\label{fig:scurve}
\end{figure}
